% !TeX program = xelatex
% Overleaf: Menu -> Compiler -> XeLaTeX
\documentclass[UTF8,11pt,a4paper]{ctexart}

\usepackage[a4paper,margin=2.45cm]{geometry}
\usepackage{amsmath,amssymb,amsthm,mathtools}
\usepackage{bm}
\usepackage{booktabs,longtable,tabularx,array}
\usepackage{enumitem}
\usepackage{xcolor}
\usepackage{microtype}
\usepackage{hyperref}

\hypersetup{
  colorlinks=true,
  linkcolor=blue!55!black,
  citecolor=green!45!black,
  urlcolor=blue!65!black,
  pdftitle={从 Weil 猜想到数域中心线：极化、过滤 Hodge 结构与黎曼猜想的存在性审计},
  pdfauthor={研究笔记整理稿}
}

\setlist{nosep,leftmargin=2em}
\renewcommand{\arraystretch}{1.16}
\setlength{\emergencystretch}{2em}

\newtheorem{theorem}{定理}[section]
\newtheorem{proposition}[theorem]{命题}
\newtheorem{lemma}[theorem]{引理}
\newtheorem{corollary}[theorem]{推论}
\theoremstyle{definition}
\newtheorem{definition}[theorem]{定义}
\newtheorem{problem}[theorem]{问题}
\theoremstyle{remark}
\newtheorem{remark}[theorem]{注}

\newcommand{\C}{\mathbb C}
\newcommand{\R}{\mathbb R}
\newcommand{\N}{\mathbb N}
\newcommand{\Tr}{\operatorname{Tr}}
\newcommand{\Spec}{\operatorname{Spec}}
\newcommand{\Ran}{\operatorname{Ran}}
\newcommand{\diag}{\operatorname{diag}}
\newcommand{\dist}{\operatorname{dist}}
\newcommand{\ord}{\operatorname{ord}}
\newcommand{\dd}{\,\mathrm d}
\newcommand{\1}{\mathbf 1}
\newcommand{\cH}{\mathcal H}
\newcommand{\cA}{\mathcal A}
\newcommand{\cM}{\mathcal M}
\newcommand{\cE}{\mathcal E}
\newcommand{\cQ}{\mathcal Q}
\newcommand{\cZ}{\mathcal Z}
\newcommand{\cC}{\mathcal C}
\newcommand{\statusT}{\textnormal{[T]}}
\newcommand{\statusU}{\statusT} % legacy TeX command name; rendered status is [T]
\newcommand{\statusC}{\textnormal{[C]}}
\newcommand{\statusE}{\textnormal{[E]}}
\newcommand{\statusN}{\textnormal{[N]}}
\newcommand{\statusR}{\textnormal{[R]}}
\newcommand{\statusO}{\textnormal{[O]}}

\title{\bfseries 从 Weil 猜想到数域中心线\\
\large 极化、过滤 Hodge 结构与黎曼猜想的存在性审计}
\author{研究笔记整理稿（第 001--202 篇）}
\date{2026 年 9 月 6 日修订（原稿：2026 年 9 月 1 日）}

\begin{document}
\maketitle

\begin{abstract}
本文从有限域上 Weil 猜想的黎曼猜想部分抽离出一个可复用的证明机制：全局迹或行列式公式把算术数据变成谱，函数方程给中心对称，正极化或有限 Hodge 负指标阻止谱向中心线外漂移。本文严格证明有限维极化 Lefschetz--Frobenius（PLF）定理、双向 tempered 动力学定理和 bounded finite-trace Hodge--Weil 蕴含；filtered primitive Weil（FPW）则被表述为带\emph{定量 divisor-mode separation} 的条件框架。本文进一步证明，在 Mellin-coherent fixed-annulus adaptive Sobolev 子类中，该 separation 可由 compact-frequency dual 自动构造；任意随尺度变化的 Gram 则仍需把它作为独立输入。

对经典 Riemann zeta 函数，本文逐项审计这些结构的存在性。显式公式、Euler--Gamma 数据、正 Gram 载体、primitive/Tate 消元、Poisson 与 Mellin 兼容、有限迹代数、Cauchy effect cone、Vaughan Type I/II 分解以及有限矩阵通道均可无条件构造。尚未得到的是一个不借用零点位置的统一算术 Hodge 估计：等价地，需控制实际算术向量的 subpower Hodge 范数，或控制 Abel Hodge current 的 Cauchy 加权负谱迹。本文进一步把后一问题无条件局部化到平方根高度核心；四分量 Vaughan Gram 又可精确商化为 canonical $2\times2$ Type I/II Gram，并可无条件删除每个高度块中 $n\leq T/\log^A T$ 的低算术长度。hard-channel Hadamard 正规形说明完整 cross-Gram 预算仍逐字等价于原 physical energy；另一方面，截断 M\"obius defect 可无条件实现为乘法 threshold complex 的有限 Hodge heat supertrace。profinite divisibility cylinders又给一个适用于 twists与 ideals的 global positive gcd polarization，由此证明在 $T\geq\log^K Y$、$K>6$ 的全部 hard arithmetic coefficient diagonals对 barrier总贡献为 $o(1)$。本文再构造 ordinary prefix discrepancy field 到 modulated current 的 exact Selberg--Volterra transport，并证明 multiscale Selberg profile 的 Hardy majorant。elementary incidence 进一步无条件消去 $N\leq T^{2-\eta}$ 的二参数楔；而在 $s\leq\log^{-K}N$ 上的 full Selberg-order profile 与 RH 等价，不是可直接调用的较弱无条件输入。未决部分因此被精确限制到 square-root resonance wedge、低 polylog height 与 continuum/Gamma joint gluing。最后证明的诸 no-free-lunch 说明：小尾部、小耦合、稳定低秩通道或 raw factorization Bessel bound 本身都不能控制物理能量，真正成功的数域 Weil 理论必须含有独立、非循环的 arithmetic Hodge 输入。

本文还证明：任意有限 ambient spectral measure 的 effects 都可完全等价地表示为长度侧的 capped positive-definite subcorrespondences；signed orbit current 的负指标是该 cone 上的精确线性最优化，不交高度 blocks 无损可加，而 prime、continuum 与 Gamma 必须先 joint gluing 再取负部。

进一步地，negative part具有 exact reciprocal-barrier变分式；任意正背景的倒数会生成新的 positive-definite correspondence kernel。由 threshold-complex harmonic currents形成的有限 predictor在该 kernel下可作 amplitude-certified Schur shorting，从而把 full quadratic energy替换为相对正 Hodge background的 one-sided residual。

本文再以 sharp square completion 删除该 pointwise amplitude 前提：加入最小二次 Hodge lift 后，背景自动保持正性，负指标由 Schur residual的 reciprocal二阶矩与 predictor的 reciprocal-cubic四阶矩控制；后者仍是显式 positive-definite correspondence Gram。平方根楔中的下一输入因而成为一条 integrated fourth-moment budget，而非 $L^\infty$ leverage。

进一步的 quartic-capacity下界表明，上述 unscaled fourth price在 zeta shell上至少为 $TD_T^2/\log T$，故它未必是较弱输入。本文于是把 reciprocal-barrier upper bound提升到不要求对易的 finite von Neumann algebra，并用 relative continuous functional calculus对 explicit predictor作 lattice clipping。所得严格正背景仅支付 one-sided quadratic clipped residual；它既不需要 pointwise leverage，也不需要 fourth moment。

非构造存在性方面，本文证明一个 Archimedean localizing completion：若每个有限 arithmetic word level 都有一致有界、渐近满足 mixed Hodge localizers 的 tracial moment functional，则 finite-intersection compactness 与 GNS 自动产生全局正表示，无需预先选择相容有限模型。与此同时，一个精确的 65 维同矩、共同范数上界反例及其 66 维同实际范数扩张表明，维数、算子范数与前四个标量迹矩即使完全相同，也不能区分正定矩阵与具有固定正比例负谱的矩阵；所以一次性的 scalar fourth-moment completion 原理上不可能推出完整 Weil 正性。

因此，本文没有证明 RH 或 GRH；它给出的是若干已证明的结构蕴含、明确标注的条件框架、存在性边界、失败机制，以及一个缩窄到具体算术不等式的下一阶段研究议程。所有内部新结果仍有待独立同行核验。
\end{abstract}

\noindent\textbf{关键词：}Weil 猜想；黎曼猜想；Hodge--Riemann 正性；显式公式；有限迹负指标；Vaughan 恒等式；Hilbert--Pólya

\medskip
\noindent\textbf{状态标签。}
\statusT{} 表示在列明假设下已有完整内部证明；
\statusR{} 表示引用的外部定理；
\statusC{} 表示实际应用依赖未证条件；
\statusE{} 表示有限数值或符号实验；
\statusN{} 表示严格反例或障碍；
\statusO{} 表示开放输入或待完成审查。
RH/GRH 等价或等强性另行写明，不用证据等级字母代替。
“无条件”只指不使用 RH 的推导，不等于已通过外部同行审查。

\tableofcontents

\section{引言：问题、方法与结论边界}

函数方程至多把一个谱点与其倒数或中心反射配对，单靠它不能推出纯性。
对于曲线，Weil 的 correspondence 论证由几何正性和 Frobenius 幂迹接口给出平方根界。
对于一般光滑射影簇，标准猜想的极化纲领是一条充分结构路线；
Deligne 的一般纯性定理不能表述为已经由本文写出的复 Hodge 内积直接证明。
这两种理论路径须分开\cite{Deligne1974,GrothendieckStandard,KleimanWeil}。

数域问题的基本困难是：每个素数处的局部 Frobenius 纯性并不控制全局 Euler 乘积解析延拓后的零点。数域对应的原始全局接口是 Weil 显式公式\cite{WeilExplicit}。即使
\[
  \zeta(s)=\prod_p(1-p^{-s})^{-1}
\]
的每个局部参数都最简单，经典 RH 仍然是开放问题。因此所需对象必须是一个\emph{全局}结构：它要同时看见所有素数幂、Gamma 因子、极点和零点，并带有不以 RH 为前提的正性或有限负指标。

本文把 173 篇文档重组为如下成果链：
\begin{enumerate}
  \item 证明有限维 PLF、幂迹、无限维 adjoint、tempered 与 tracial determinant 等结构定理；
  \item 构造带定量对偶分离公理的 filtered primitive Weil package，把“纯权”条件化为算术向量随尺度的 subpower 增长；
  \item 对 zeta 无条件实现解析、Euler、positive-carrier 与有限化部分，并在 adaptive Sobolev carrier 上建立定量分离；最后的 arithmetic tightness 仍与 RH 等强；
  \item 从 Nyman--Beurling、prime graph、wavelet、Green kernel、passivity 等多条路线提取等价或充分判据，同时记录不能闭合的 no-go；
  \item 把最新路线压缩为有限迹负指标，再局部化到平方根核心、canonical 二通道 Vaughan quotient 与 $n>T/\log^A T$ 的 hard arithmetic lengths；
  \item 证明 hard cross-Gram 目标的 Hadamard 循环性与 raw factorization multiplicity 障碍，并把截断 M\"obius defect 实现为 threshold-complex Hodge heat supertrace；
  \item 构造 profinite divisibility polarization，将 M\"obius second moment转成任意单调权，并无条件消去高于 polylog height 的 hard coefficient diagonals；
  \item 构造 Selberg--Volterra cross-fiber transport，把 modulated near-products严格归约为 ordinary multiscale prefix profile 的 Hardy cost；
  \item 证明 elementary prefix incidence 无条件消去 $N\leq T^{2-\eta}$ 的 sub-square-root wedge，并以 telescoping--Mellin argument证明 polylogarithmic full profile 与 RH 等价；
  \item 把频谱 effect cone 精确搬到长度侧的 capped correspondence kernels，证明其负指标对偶、dyadic block gluing 与 joint-orbit cancellation，并给出纯 kernel 的 bounded-index Weil 定理；
  \item 证明 exact reciprocal-barrier identity，以 $d\mu/B$ 构造正定 kernel，并把 threshold-complex predictor的贡献精确化为带 amplitude 证书的 Schur residual；
  \item 以最优二次 square completion 自动保持背景正性，把 amplitude 证书替换为 Schur residual 与显式四阶 correspondence Gram；
  \item 证明四阶代价的 capacity必要下界，并以 finite-trace relative lattice clipping进一步把它替换为 one-sided quadratic residual；
  \item 证明仅靠正 Gram 的低秩压缩不可能“免费”产生 Hodge--Riemann 约束，因而明确下一步必须构造何种新结构。
  \item 证明 scalar fourth-moment completion 的 65 维严格 no-go，并给出增长 mixed localizers 的 Archimedean 非构造补全定理。
\end{enumerate}

本文区分已证明蕴含与条件框架：PLF、tempered、有限迹 effect 对偶、bounded-defect normality 和 no-free-lunch 属于 \statusT{}；FPW 中心线定理在定量模式分离公理下属于 \statusC{}；实际 zeta tightness 或 bounded negative trace 属于 \statusO{}，其 RH 强度须另行审计。有限浮点实验统一标为 \statusE{}，只用于发现结构、排除错误目标，不构成 RH 证据。

\section{有限域原型：极化 Lefschetz--Frobenius 结构}

\subsection{最小谱论机制}

\begin{theorem}[\statusT{} 正 similitude 强制纯性]
设 $H$ 是有限维复向量空间，$F\in\mathrm{GL}(H)$，$Q>0$。若存在正定 Hermite 内积 $h$ 使
\[
  h(Fx,Fy)=Qh(x,y),
\]
则 $Q^{-1/2}F$ 为酉算子；因而 $F$ 可对角化，且每个特征值 $\alpha$ 满足 $|\alpha|=Q^{1/2}$。
\end{theorem}

\begin{proof}
令 $U=Q^{-1/2}F$，则 $h(Ux,Uy)=h(x,y)$。有限维酉谱定理给出结论。
\end{proof}

这个定理本身过于抽象：若已经知道谱在圆周上，常可反向制造内积。有效性来自极化必须由独立的代数与几何结构产生。

\subsection{PLF 代数与广义有限维 Weil 定理}

\begin{definition}[PLF 代数]
固定 $q>1$ 与整数 $d\geq0$。一个极化 Lefschetz--Frobenius 代数由下列数据组成：
\begin{enumerate}
  \item 有限维分次交换复代数 $H=\bigoplus_{n=0}^{2d}H^n$，
  对齐次元素满足 $xy=(-1)^{\deg x\deg y}yx$，带保分次的反线性代数对合 $x\mapsto\bar x$；
  \item 顶次迹 $\tau:H^{2d}\to\C$，使乘法配对 $H^n\times H^{2d-n}\to\C$ 非退化；
  \item 实 Lefschetz 元 $L\in H^2$，且 $L^{d-n}:H^n\to H^{2d-n}$ 在 $n\leq d$ 时为同构；
  \item 分次代数自同构 $F$，满足 $FL=qLF$、$F\bar x=\overline{Fx}$ 和 $\tau(Fz)=q^d\tau(z)$；
  \item 对 $0\leq a\leq d$ 定义 primitive 部分
  \[
    P^a=\ker\!\left(L^{d-a+1}:H^a\longrightarrow H^{2d-a+2}\right),
  \]
  并要求 Lefschetz 直和分解 $H^n=\bigoplus_{a+2r=n}L^rP^a$，
  其中 $0\leq a\leq d$、$0\leq r\leq d-a$，范围外分次约定为零；
  \item 对每个 $a$ 另给一个可逆复线性 primitive 极化算子
  $J_a:P^a\to P^a$，与实结构及 $F|_{P^a}$ 交换。
  对每个允许的 $(a,r)$ 选取非零实常数 $c_{a,r}$，逐 primitive 分量定义
  \[
    S_n(L^ru)=c_{a,r}L^{d-a-r}J_au,
    \qquad u\in P^a,\quad n=a+2r;
  \]
  再选取 $|\varepsilon_n|=1$，并把
  \[
    h_n(x,y)=\varepsilon_n\tau\bigl(xS_n(\bar y)\bigr)
  \]
  是 Hermite 对称且正定作为 Hodge--Riemann 极化公理。
\end{enumerate}
\end{definition}

\begin{remark}[奇次极化是额外输入]\label{rem:odd-polarization}
旧定义相当于所有 $J_a=I$。若存在非零奇次 primitive 实类 $u$，
则分次交换性给 $u^2=0$，从而
$h_a(u,u)=\varepsilon_ac_{a,0}\tau(L^{d-a}u^2)=0$，与正定性矛盾。
尤其通常的正亏格曲线 $H^1$ 不能实例化旧定义；统一复相位不能修复它。
修订后必须独立构造 $J_a$ 并证明正性和 Frobenius 相容性，
不能从已知纯谱反造后声称得到几何证明。曲线的 degree/Rosati 对应正性
基准是另一种结构，不能自动称为这个 PLF 定义的实例。
\end{remark}

\begin{theorem}[\statusT{} 有限维广义 Weil 结构定理]\label{thm:plf}
对任意 PLF 代数，$F|_{H^n}$ 可对角化，且
\[
  \Spec(F|_{H^n})\subset\{\alpha\in\C:|\alpha|=q^{n/2}\}.
\]
\end{theorem}

\begin{proof}
若 $x=L^ru$、$u$ primitive 且 $n=a+2r$，由 $FL=qLF$ 及 $J_aF=FJ_a$ 得到
\[
  S_nF=q^{n-d}FS_n.
\]
再用 $F$ 保乘法、与共轭交换以及 $\tau(Fz)=q^d\tau(z)$，可得
\[
  h_n(Fx,Fy)=q^nh_n(x,y).
\]
应用上一最小谱论定理，取 $Q=q^n$ 即得。
\end{proof}

\begin{theorem}[\statusT{} 迹公式到权重线]
若另有点数恒等式
\[
  N_m=\sum_n(-1)^n\Tr(F^m|H^n),\qquad
  Z(T)=\exp\!\left(\sum_{m\geq1}\frac{N_mT^m}{m}\right),
\]
则
\[
  Z(T)=\prod_n\det(1-TF|H^n)^{(-1)^{n+1}}.
\]
第 $n$ 次上同调贡献的 reciprocal 零点或极点满足 $|\alpha|=q^{n/2}$；令 $T=q^{-s}$，其 divisor 位于 $\Re s=n/2$。
\end{theorem}

\begin{proof}
逐次使用 $-\log\det(1-TF)=\sum_{m\geq1}\Tr(F^m)T^m/m$，再应用定理~\ref{thm:plf}。
\end{proof}

\subsection{correspondence 版本}

\begin{theorem}[\statusT{} reciprocal pairing 与幂迹界]
设非零复数多重集 $\{\alpha_j\}_{j=1}^N$ 在 $\alpha\mapsto Q/\alpha$ 下不变，并且
\[
  \left|\sum_j\alpha_j^m\right|\leq C Q^{m/2}\qquad(m\geq1).
\]
则每个 $|\alpha_j|=Q^{1/2}$。
\end{theorem}

\begin{proof}
令 $\beta_j=\alpha_j/Q^{1/2}$。幂迹界说明 $\sum_{m\geq1}(\sum_j\beta_j^m)z^m$ 的收敛半径至少为 $1$；而该级数等于 $\sum_j\beta_jz/(1-\beta_jz)$，其最近极点给 $\max|\beta_j|\leq1$。倒数配对再给 $\min|\beta_j|\geq1$。
\end{proof}

这正抽象了 Weil 对曲线的原始论证：Hodge 指数或 Rosati 正性不必先成为整个复上同调的内积，只要它能对全部 Frobenius 幂产生统一平方根界。

\section{无限维、tempered 与 tracial 推广}

\subsection{Hilbert--Pólya 型 adjoint 关系}

\begin{theorem}[\statusC{} 极化生成元定理]\label{thm:adjoint}
令 $\Theta$ 是 Hilbert 空间上的稠密定义闭算子。若
\[
  D(\Theta)=D(\Theta^*),\qquad \Theta^*=cI-\Theta,
\]
则 $A=\Theta-cI/2$ 反自伴，且 $\Spec(\Theta)\subset c/2+i\R$。若一个不借用零点位置构造的正规化行列式满足
\[
  \Lambda(s)=E(s)\det_\infty(s-\Theta)
\]
且 $E$ 的 divisor 已知，则 $\Lambda$ 的相应零点位于中心线。
\end{theorem}

\begin{proof}
$A^*=-A$。反自伴算子的谱位于虚轴；对特征向量也可直接由
$2\Re\rho\|v\|^2=\langle(\Theta+\Theta^*)v,v\rangle=c\|v\|^2$ 得到。
\end{proof}

\begin{remark}[循环性]
以已知零点为正交基并令 $\Theta e_\rho=\rho e_\rho$，会把 RH 偷放进结构中。数域存在性定理必须同时说明：对象由算术数据构造；正性不使用 RH；迹/行列式恢复完整 Gamma--Euler divisor；无限维定义域和正规化可控。
\end{remark}

\subsection{双向 tempered 动力学}

\begin{theorem}[\statusT{} 有限维 tempered 纯性与酉化]\label{thm:tempered}
令 $F$ 是有限维 Hermite 空间上的可逆算子，并置 $U=q^{-w/2}F$。若对每个 $\epsilon>0$ 有
\[
  \|U^n\|\leq C_\epsilon e^{\epsilon|n|}\qquad(n\in\mathbb Z),
\]
则 $\Spec(F)$ 位于 $|\alpha|=q^{w/2}$。若进一步 $\sup_{n\in\mathbb Z}\|U^n\|\leq M$，则存在与原内积等价的正内积使 $U$ 精确酉。
\end{theorem}

\begin{proof}
谱半径公式分别作用于 $U$ 与 $U^{-1}$，给每个谱值同时满足 $|\lambda|\leq1$ 与 $|\lambda|\geq1$。在一致有界情形，Ces\`aro 度量
\[
  H_N=\frac1{2N+1}\sum_{n=-N}^N(U^n)^*U^n
\]
满足 $M^{-2}I\leq H_N\leq M^2I$。有限维有界闭集紧，故可取范数收敛子列；边界项为 $O(N^{-1})$，极限 $H_\infty$ 满足 $U^*H_\infty U=H_\infty$。
\end{proof}

一般无限维 Hilbert 空间中不能使用算子范数紧性。对一致有界的强连续群 $G_t$，应取
\[
 H_T=\frac1{2T}\int_{-T}^{T}G_t^*G_t\,\dd t
\]
在 weak-operator compact 单位球中的 cluster subnet；区间平移只产生 $O(T^{-1})$ 的边界，故极限满足 $G_s^*H_\infty G_s=H_\infty$。双向 subexponential 增长则由 resolvent 的 Laplace 表示迫使生成元谱位于 $i\R$。所得 tempered Weil 对象对直和、张量积、对偶、子商、Schur 函子和 Tate twist 封闭。这里的张量结论只适用于已经具有\emph{全局}谱包的对象，不能从局部 Satake 参数的张量积直接推出 Rankin--Selberg 全局零点空间。

\subsection{迹维数与零点重数}

令 $\Gamma$ 为离散谱集，$m_\gamma\in\N$。在 $\ell^\infty(\Gamma)$ 上定义
\[
  \tau(a)=\sum_{\gamma\in\Gamma}m_\gamma a(\gamma),\qquad a\geq0.
\]
这是 faithful normal semifinite trace，且坐标投影满足 $\tau(P_\gamma)=m_\gamma$。

这里使用的 determinant 不是 Fuglede--Kadison determinant。设谱集局部有限，令
\[
 \nu_\rho=\tau_{\rm gr}(P_\rho)\in\mathbb Z,
 \qquad
 E_g(z)=(1-z)\exp\!\left(\sum_{k=1}^{g}\frac{z^k}{k}\right),
\]
并假设某个整数 $g\geq0$ 使
$\sum_{\rho\ne0}|\nu_\rho||\rho|^{-g-1}<\infty$。固定次数不超过 $g$ 的正规化多项式 $P$ 后，定义带 tracial exponents 的 Weierstrass 谱乘积
\[
 \det^{\mathrm W}_{\tau,g}(s-\Theta)
 =e^{P(s)}s^{\nu_0}\prod_{\rho\ne0}E_g(s/\rho)^{\nu_\rho}.
\]
其对数导数是 regularized resolvent trace：
\[
 \frac{\dd}{\dd s}\log\det^{\mathrm W}_{\tau,g}(s-\Theta)
 =P'(s)+\frac{\nu_0}{s}
  +\sum_{\rho\ne0}\nu_\rho
   \left(\frac1{s-\rho}+\sum_{k=0}^{g-1}\frac{s^k}{\rho^{k+1}}\right).
\]
因此正规化和 polynomial counterterm 不改变每个谱点的局部 order $\nu_\rho$。

\begin{theorem}[\statusC{} tracial polarized Weil spectral product]
设 $(\cM,\tau)$ 是 semifinite tracial von Neumann 代数，$\Theta$ 是 affiliated normal 算子，谱投影局部有限且具有有限整数 graded trace，并满足上述 genus-$g$ summability。若
\[
  \Theta^*=c-\Theta,\qquad
  \Lambda(s)=E(s)\det^{\mathrm W}_{\tau,g}(s-\Theta),
\]
其中 $E$ 的 divisor 已知，则非平凡 divisor 位于 $\Re s=c/2$，且谱乘积以 $\nu_\rho=\tau_{\rm gr}(P_\rho)$ 恢复有符号代数重数。
\end{theorem}

\begin{proof}
中心线结论来自定理~\ref{thm:adjoint}。上述 canonical product 的 $\rho$ 因子具有局部 order $\nu_\rho$；其余因子在该点全纯且非零。已知因子 $E$ 再把结论转移到 $\Lambda$。
\end{proof}

这说明普通 Hilbert 空间本征子空间维数不是恢复重数的唯一方式。标量最小 GNS fiber 可以是一维，而 trace dimension 仍精确携带重数；但必须把这种全纯谱乘积与正实值的 Fuglede--Kadison determinant 明确区分。

\subsection{非构造 localizer 补全与低阶矩边界}

下面回答一个自然的存在性问题：是否可以不显式构造数域 Hilbert 空间，而由
所有有限 arithmetic configurations 的可满足性非构造地产生全局 Weil 型正表示？
答案是肯定的，但必须使用随 word degree 增长的 matrix localizers；固定的前四个
标量矩严格不够。

\begin{theorem}[\statusN{} scalar fourth-moment completion no-go]
\label{thm:four-moment-nogo}
存在两个 $65\times65$ Hermitian 矩阵 $H_+,H_-$，满足
\[
 \|H_+\|=1,\quad\|H_-\|=4/5,\quad
 \frac1{65}\Tr H_+^k=\frac1{65}\Tr H_-^k
 \quad(0\le k\le4),
\]
但 $H_+\succ0$，而 $H_-$ 有一个负特征值。共同的矩依次为
\[
 1,\quad \frac7{13},\quad \frac{109}{325},\quad
 \frac{371}{1625},\quad \frac{1357}{8125}.
\]
任取多个直和后，负谱比例仍为 $1/65$。两者具有共同范数上界 $1$。
若要求实际范数完全相同，令 $\widehat H_\pm=H_\pm\oplus[1]$：
它们维数为 $66$、范数同为 $1$、前四矩仍相同，
归一化矩变为 $(65m_k+1)/66$，负谱比例为 $1/66$。
因此任何只依赖维数、统一算子范数和
$m_0,\ldots,m_4$ 的抽象补全原则，都不能推出 Hodge current 的完整正性。
\end{theorem}

\begin{proof}
取归一化谱测度
\[
 \mu_+=\frac3{13}\delta_{1/5}
 +\frac9{13}\delta_{3/5}+\frac1{13}\delta_1,
\]
\[
 \mu_-=\frac1{65}\delta_{-1/5}
 +\frac8{13}\delta_{2/5}+\frac{24}{65}\delta_{4/5}.
\]
相应重数是 $(15,45,5)$ 与 $(1,40,24)$。对 $0\le k\le4$，恒等式
\[
 -\left(-\frac15\right)^k+15\left(\frac15\right)^k
 -40\left(\frac25\right)^k+45\left(\frac35\right)^k
 -24\left(\frac45\right)^k+5=0
\]
给出矩相等。正负谱结论从支撑直接可见。
\end{proof}

甚至一阶 localizer
\[
 \begin{pmatrix}m_1&m_2\\m_2&m_3\end{pmatrix}
\]
对二者相同且行列式为 $1104/105625>0$。但增长一级已经可以检测负方向：
对 $p(x)=(x-2/5)(x-4/5)$，
\[
 \int xp(x)^2\,d\mu_-(x)=-\frac9{8125}<0.
\]
这说明需要增长的 localizing word spaces，或额外的 flatness/递推关系。

\begin{theorem}[\statusT{} Archimedean finite-satisfiability completion]
\label{thm:localizer-completion}
令 $\cA_0$ 为具有可数 word 张成集的幺正复 $*$-代数，$Q\subset(\cA_0)_h$ 为
Archimedean quadratic module：对每个 $a\in\cA_0$，存在 $R_a<\infty$ 使
\[
 R_a^2-a^*a\in Q.
\]
令 $h=h^*\in\cA_0$，并在一个可数 arithmetic word family $E$ 上预定
length-side moments $\lambda(e)$。

假设对每个有限约束集 $F$ 和每个 $n\ge1$，存在定义在包含 $F$ 所涉 words
的有限维 $*$-空间上的线性泛函 $L_{F,n}$，满足：
\begin{enumerate}
 \item $L_{F,n}(1)=1$，所选 $*$-关系与迹关系的误差不超过 $1/n$；
 \item 所选 Gram 与 $Q$-localizing matrices 半正定到误差 $1/n$；
 \item 对所选 words $u_1,\ldots,u_r$，
 \[
 [L_{F,n}(u_i^*hu_j)]_{i,j}
 \succeq-\frac1n[L_{F,n}(u_i^*u_j)]_{i,j};
 \]
 \item $|L_{F,n}(e)-\lambda(e)|\le1/n$ 对所选 $e\in E$ 成立。
\end{enumerate}
则存在全局 tracial state $L$ 与其 GNS 表示 $(\pi,\cH,\Omega)$，使
\[
 L(e)=\lambda(e)\quad(e\in E),\qquad \pi(h)\succeq0.
\]
\end{theorem}

\begin{proof}
固定一个可数线性基。对任意有限基坐标集 $K$，加入每个坐标 $a$ 的 $\{1,a\}$ Gram、
Archimedean 元 $R_a^2-a^*a$ 的标量约束及所涉 words。
矩阵不等式均取 Hermitian part；先将有限泛函 Hermitian 化为
$\widetilde L(x)=(L(x)+\overline{L(x^*)})/2$。对 $\epsilon=1/n\le1$，
\[
|\widetilde L(a)|^2\le(1+\epsilon)(\widetilde L(a^*a)+\epsilon)
\le(1+\epsilon)(R_a^2+2\epsilon)\le2(R_a^2+2).
\]
$Q$ 未必可数，故不使用可数约束穷尽。按 $(K,F,n)$ 构成有向网：
$F$ 是任意有限约束集，$K$ 包含其全部基坐标，按包含关系及 $n$ 的大小排序。
有限可满足性另纳入 $K$ 的上述坐标界及辅助 words；只记录 $K$ 的坐标，
其他坐标补 $0$。这些点在固定紧圆盘乘积中，紧致性给收敛子网。
每条固定约束最终均被包含、误差趋零，且只依赖有限坐标并闭，
故极限是满足全部精确约束的 tracial state $L$。
精确界 $|L(a)|\le R_a$ 只在此时使用。
GNS 构造中
\[
 \langle\pi(h)\pi(p)\Omega,\pi(p)\Omega\rangle=L(p^*hp)\ge0.
\]
由 $p^*(R_a^2-a^*a)p\in Q$ 得 $L(p^*a^*ap)\le R_a^2L(p^*p)$，故各 $\pi(a)$ 有界，故 $\pi(h)\succeq0$。
\end{proof}

定理~\ref{thm:localizer-completion} 与 truncated tracial moment 和
flat-extension 理论的方向一致\cite{BurgdorfKlep,MourrainSchmudgen}，但不假设
低阶 flatness。它的价值是消除“如何相容地选择所有有限模型”这一选择问题；
它不免费证明局部正性。若全局交为空，紧致性的反面给某个有限 word level 已
不可满足；有限维凸分离给分离证据。可计算的 SDP/SOS 对偶证书仍需独立的编码、对偶性及精度条件。

对 zeta 的非循环应用还必须验证：$\cA_0,Q,h$ 完全由 length side 定义；
bounded resolvent 或 Cayley-transform moments 在极限中恢复 divisor visibility；
并且相应 tracial Weierstrass product 确实恢复 Gamma--Euler divisor。若把全部
Weil tests 的正性直接列入有限约束，前提仍与 RH 等价。真正可能较弱的目标是
找到一个小而代数闭合的 arithmetic operator system，使其 mixed localizers
可由 Type I/II、Gamma 与 continuum identities 控制，而闭包自动看见所有负方向。

\section{Filtered primitive Weil 条件框架}

精确酉性对数域过强。尺度 $X$ 上的过滤模型只需辨认每个固定 divisor mode 的增长指数，并证明实际算术向量的 Hodge 范数没有幂次增长。

\begin{definition}[定量 FPW package]
中心为 $c/2$ 的 quantitative filtered primitive Weil package 包含：
\begin{enumerate}
  \item 由有限 Euler 数据构造的 arithmetic vector $v_X$，先减去 continuum/Tate 主模态；
  \item 正半定 Gram/Hodge form $\cQ_X$；任何替换必须区分 $X^{o(1)}$-conditioned quotient 同构与仅对 actual vector 成立的 additive $X^{o(1)}$ enclosure，后者可转移 tightness 结论但不自动转移全空间 dual；
  \item primitive dilation polynomial $P_X(U)$，在每个固定临界谱紧集上以 $X^{o(1)}$ 双侧可逆；
  \item 与显式公式相容的定性 divisor expansion：每个固定 $\rho$ 的形式模态为
  \[
    X^{\rho-c/2}m_X(\rho),\qquad |m_X(\rho)|=X^{o(1)}
  \]
  （双侧意义），且 off-center divisor 不对全部 fibers 恒不可见；
  \item \emph{定量模式分离}：在商空间 $V_X/\ker\cQ_X$ 上，对每个固定 divisor point $\rho$ 存在对偶泛函 $\ell_{X,\rho}$ 和一条共尾子列，使
  \[
   \|\ell_{X,\rho}\|_{\cQ_X^*}=X^{o(1)},\qquad
   |\ell_{X,\rho}(v_X)|
      \geq X^{\Re\rho-c/2-o(1)}.
  \]
  这里
  $\|\ell\|_{\cQ_X^*}=\sup_{\cQ_X(v)>0}|\ell(v)|/\cQ_X(v)^{1/2}$；特别要求 $\ell$ 在 $\ker\cQ_X$ 上消失；
  \item 高频、边界、Taylor、sampling 与 finite-completion tails 的总范数为 $X^{o(1)}$，且有限阶显式公式给
  \[
    \cQ_X(v_X)\leq X^{\max(0,2\Theta_+-c)+o(1)};
  \]
  \item arithmetic tightness：$\cQ_X(v_X)=X^{o(1)}$。
\end{enumerate}
第 5 项严格强于“模态不恒为零”的定性 Mellin 唯一性；它是把 divisor coefficient 转成 Hodge 范数下界所必需的 Riesz/frame 型输入。对一般 zeta Gram realization，该项仍需单独建立；但下述定理会在 Mellin-coherent adaptive Sobolev carrier 上无条件构造它。第 7 项则是数域 Hodge--Riemann 对实际 arithmetic cycle 的内容。
\end{definition}

\begin{theorem}[\statusC{} quantitative filtered primitive Weil 中心线定理]\label{thm:fpw}
若非空、中心对称的有限阶 divisor data 具有上述 quantitative FPW package，且
$\Theta_+=\sup_\rho\Re\rho<\infty$，则
\[
  \limsup_{X\to\infty}
  \frac{\log\max(1,\cQ_X(v_X))}{\log X}
  =\max(0,2\Theta_+-c).
\]
特别地，arithmetic tightness 推出全部 divisor 位于 $\Re\rho=c/2$。
\end{theorem}

\begin{proof}
第 6 项给所需上界。固定 $\rho$，对偶 Cauchy--Schwarz 严格给出
\[
 |\ell_{X,\rho}(v_X)|^2
 \leq \|\ell_{X,\rho}\|_{\cQ_X^*}^2\cQ_X(v_X),
\]
故沿第 5 项的共尾子列
\[
  \cQ_X(v_X)\geq X^{2\Re\rho-c-o(1)}.
\]
对任意 $\varepsilon>0$ 选取 $\Re\rho>\Theta_+-\varepsilon$，再令 $\varepsilon\downarrow0$，得到 limsup 的反向不等式。中心对称和 divisor 非空给 $\Theta_+\geq c/2$。tightness 迫使 $\Theta_+\leq c/2$；中心反射对称再排除左侧零点。
\end{proof}

\begin{remark}[审计后的降级]
若删除第 5 项，任意 scale-varying Gram 的定性可见性不能排除 subpower 精度的相消，因而不足以证明上述 exponent identity。弱版本应视为 \statusO{} 研究框架。下面的正面结果依赖 fixed-annulus fibers 来自同一个 Dirichlet series 的跨尺度 Mellin coherence，不能反推任意 Gram 自动满足第 5 项。
\end{remark}

\begin{theorem}[\statusT{} Mellin-coherent Sobolev 定量分离]\label{thm:mellin-dual}
固定 $A_0>1$、$\kappa>0$，令 $B=\log A_0$。设
\[
 L_a(s)=\sum_{n\geq1}a(n)n^{-s}
\]
在某个右半平面绝对收敛，并亚纯延拓到目标 divisor point $\rho$，且在 $\rho$ 有真正极点。对支撑于 $X/A_0<n\leq A_0X$ 的有限向量 $x$ 定义
\[
 F_{X,x}(\tau)=\sum_nx_n n^{-c/2}e^{-i\tau\log(n/X)},\qquad
 Q_{X,r}(x)=\frac1{2\pi}\int_\R
 \frac{|F_{X,x}(\tau)|^2}{(\tau^2+\kappa^2)^r}\,\dd\tau .
\]
这里 $r\in\mathbb N_{\ge1}$，有限指数和有界且 $(\tau^2+\kappa^2)^{-r}$ 可积，
故 $Q_{X,r}$ 有限。取实际向量 $x_{X,n}=a(n)\1_{(X/A_0,A_0X]}(n)$。
若 $r_X\in\mathbb N_{\ge1}$ 且 $r_X=o(\log X)$，则对每个固定 $\rho$ 存在 compact-frequency 对偶泛函 $\ell_{X,\rho}$ 和 $X_j\to\infty$，使
\[
 \|\ell_{X_j,\rho}\|_{Q_{X_j,r_{X_j}}^*}=X_j^{o(1)},\qquad
 |\ell_{X_j,\rho}(x_{X_j})|
 \geq X_j^{\Re\rho-c/2-o(1)}.
\]
从而
\[
 Q_{X_j,r_{X_j}}(x_{X_j})
 \geq X_j^{2\Re\rho-c-o(1)}.
\]
\end{theorem}

\begin{proof}
置 $z_\rho=\rho-c/2$。取有界 $\chi_0\geq0$，支撑于固定紧集且 $\int\chi_0>0$，并令
\[
 h_\rho(\tau)=\int_{-B}^{B}e^{(z_\rho-i\tau)u}\,\dd u,
 \qquad q_\rho(\tau)=\chi_0(\tau)h_\rho(\tau),
\]
\[
 \ell_{X,\rho}(x)=\frac1{2\pi}\int_\R
 F_{X,x}(\tau)\overline{q_\rho(\tau)}\,\dd\tau .
\]
若 $\operatorname{supp}q_\rho\subset[-T,T]$，带权 Cauchy--Schwarz 给
\[
 \|\ell_{X,\rho}\|_{Q_{X,r}^*}^2
 \leq \frac1{2\pi}\int |q_\rho(\tau)|^2
       (\tau^2+\kappa^2)^r\,\dd\tau
 \leq C_\rho R_\rho^{2r}.
\]
所以 $r_X=o(\log X)$ 时 dual norm 为 $X^{o(1)}$。

写
\[
 k_\rho(u)=\frac1{2\pi}\int
 \overline{q_\rho(\tau)}e^{-i\tau u}\,\dd\tau,\qquad
 K_\rho(z)=\int_{-B}^{B}k_\rho(u)e^{zu}\,\dd u .
\]
逐项令 $u=\log(n/X)$，在初始收敛半平面得到
\[
 \int_0^\infty \ell_{X,\rho}(x_X)X^{-z-1}\,\dd X
   =K_\rho(z)L_a(c/2+z).
\]
而 Fubini 给
\[
 K_\rho(z_\rho)=\frac1{2\pi}\int
 \chi_0(\tau)|h_\rho(\tau)|^2\,\dd\tau>0.
\]
故右端在 $z_\rho$ 保留真正极点。若左端取值的 limsup 指数小于 $\Re z_\rho$，其 Mellin 积分会在包含 $z_\rho$ 的更大右半平面绝对收敛并全纯，和该极点矛盾。因此取值沿某子列满足所写下界；再用前述 dual Cauchy--Schwarz 即得 Hodge 下界。
\end{proof}

\begin{corollary}[\statusT{} zeta adaptive Sobolev FPW4b]\label{cor:zeta-fpw4b}
对经典 zeta 取 $c=1$、$A_0=4$、$\kappa=1/2$ 和 $a(n)=\Lambda(n)-1$。此时
\[
 L_a(s)=-\frac{\zeta'}{\zeta}(s)-\zeta(s).
\]
$s=1$ 的两个主极点相消，而每个非平凡零点 $\rho$ 仍给 residue $-m_\rho\ne0$。因此任意满足 $r_X\in\mathbb N_{\ge1}$、
$r_X=o(\log X)$ 的 adaptive Sobolev carrier 无条件满足 FPW4b。sampled 与 logarithmic-moment actual energies 可由已证的 subpower additive approximation 继承中心线蕴含；该事实本身不声称它们各自已有全空间 Riesz lower bound。
\end{corollary}

\begin{proposition}[\statusT{} finite Schur dual certificate]
若有限 fiber 上 $Q(x)=x^*Gx$、$G>0$，且 $\ell(x)=d^*x$，则
\[
 \|\ell\|_{Q^*}^2=d^*G^{-1}d,
\]
并且
\[
 \begin{pmatrix}G&d\\d^*&\eta\end{pmatrix}\geq0
 \quad\Longleftrightarrow\quad
 \eta\geq d^*G^{-1}d.
\]
定理~\ref{thm:mellin-dual} 中的 weighted Cauchy--Schwarz 给出显式选择 $\eta=C_\rho R_\rho^{2r}$。
\end{proposition}

\begin{proof}
第一式是有限维 Riesz representation；第二式是对正定块 $G$ 的 Schur complement。
\end{proof}

若 Gram family 对平移协变且一致 tight，则可作 GNS completion，得到强连续酉群和 $\Theta=c/2+iA$；因此 strong FPW 会恢复真正的 Hilbert--P\'olya 型算子。filtered FPW 只控制 distinguished arithmetic cyclic vector，足以推出 divisor purity，却不自动给整个算子上的半单性。

\section{数域主定理：有限迹 Hodge 负指标}

\subsection{effect cone 上的精确对偶}

\begin{theorem}[\statusT{} 有限迹 Hodge 指数对偶]\label{thm:trace-duality}
令 $(\cM,\tau)$ 是有限 von Neumann 代数，$\tau(1)=1$，$H=H^*$ 可积。写 $H=H_+-H_-$。则
\[
  \inf_{0\leq E\leq1}\tau(HE)=-\tau(H_-),
\]
且极小值由负谱投影 $E=\1_{(-\infty,0)}(H)$ 达到。
\end{theorem}

\begin{proof}
对任意 effect，
\[
 \tau(HE)=\tau(H_+^{1/2}EH_+^{1/2})-
 \tau(H_-^{1/2}EH_-^{1/2})\geq-\tau(H_-).
\]
取 $H_-$ 的支撑投影即取等。
\end{proof}

这一定理说明正确的数域“负指标”不是最小特征值，而是所有负方向按 ambient spectral capacity 加权后的总质量。很深但很窄的 resonance 井可以破坏 operator-norm 下界，却仍具有有界迹负指标。

\subsection{bounded-defect normality}

记右半平面为 $\C_+=\{z:\Re z>0\}$。设 $F_n$ 在 $\C_+$ 全纯，$\delta_n\downarrow0$，并令
\[
 W_n(t)=[-\Re F_n(\delta_n+it)]_+,
 \qquad
 J_n=\int_\R\frac{W_n(t)}{1+t^2}\dd t<\infty.
\]
所谓 \emph{Poisson admissibility} 是指对每个 $x>\delta_n$、$y\in\R$ 都有
\[
 \Re F_n(x+iy)\geq-\frac1\pi\int_\R
 \frac{(x-\delta_n)W_n(t)}{(x-\delta_n)^2+(y-t)^2}\dd t.
\]
这一定义同时固定了符号、归一化和边界可积条件。

\begin{lemma}[\statusT{} 单点 Carath\'eodory 正规族]
若全纯函数族在每个紧集上实部有统一下界，且在一个点取值有界，则该族局部有界，因而为正规族。
\end{lemma}

\begin{proof}
在圆盘上对 $f+A$ 使用 positive-real Carath\'eodory 估计；再由重叠圆盘链把基点界传播到任意紧集，最后用 Montel 定理。
\end{proof}

\begin{theorem}[\statusC{} bounded finite-trace Hodge--Weil 主定理]\label{thm:main}
设 $\Phi$ 是 real-type、关于中心轴自对偶、阶至多一的 entire divisor。假设：
\begin{enumerate}
 \item $F_n$ 在 $\C_+$ 全纯、$\delta_n\downarrow0$，上述 $W_n,J_n$ 有定义且满足完整 Poisson admissibility 不等式；
 \item 存在非空开集 $U\subset\C_+$，使 $F_n\to\Phi'/\Phi$ 在 $U$ 的紧集上一致；
 \item 存在规范化有限迹代数 $(\cM_n,\tau_n)$、$\tau_n$-可积自伴 Hodge current $H_n$ 和误差 $\eta_n\geq0$，并已验证
 \[
   \frac{J_n}{\pi}\leq\tau_n((H_n)_-)+\eta_n;
 \]
 \item $\sup_n\{\tau_n((H_n)_-)+\eta_n\}<\infty$。
\end{enumerate}
则 $\Phi$ 的全部非零零点位于中心轴。
\end{theorem}

\begin{proof}
定理~\ref{thm:trace-duality} 把 effect cone 上的最坏 Hodge 深度精确识别为 $\tau_n((H_n)_-)$，所以 $\sup_nJ_n<\infty$。Poisson kernel在任一右半平面紧集上由 $C_K/(1+t^2)$ 控制，故 $\Re F_n$ 有统一下界。Euler 开集提供一个有界基点，正规族引理给全右半平面的局部有界性。任意子列极限在 Euler 开集等于 $\Phi'/\Phi$，由恒等定理唯一，故得到其全右半平面全纯延拓。若 $\Phi$ 在右半平面有零，则 $\Phi'/\Phi$ 在那里有不可去极点，矛盾。自对偶把左侧零点反射到右侧，于是零点全在中心轴。
\end{proof}

定理~\ref{thm:main} 是本文最宽的“广义结构定理”。它允许非交换代数、无限秩负谱、趋于 $-\infty$ 的最小特征值；只要负谱的规范迹质量不发散，中心线仍成立。有限域 Hodge--Riemann 正性对应更强的 $H_n\geq0$。

\section{经典 zeta 的规范有限迹实现}

对 $z\in\C_+$ 置 $s=1/2+z$，并定义只使用 prime powers、Gamma 数据和连续主密度的 Abel candidate
\begin{align*}
 S_Y(s)&=\sum_{n\geq2}\Lambda(n)n^{-s}e^{-n/Y},\\
 I_Y(s)&=\int_1^\infty x^{-s}e^{-x/Y}\dd x
       =Y^{1-s}\Gamma(1-s,1/Y),\\
 A_\infty(s)&=\frac1s-\frac12\log\pi
              +\frac12\frac{\Gamma'}{\Gamma}(s/2),\\
 F_Y(z)&=A_\infty(s)+I_Y(s)-S_Y(s).
\end{align*}
指数权保证 $S_Y,I_Y$ 为 entire functions，因而每个 $F_Y$ 在 $\C_+$ 全纯；在 $\Re s>1$，它局部一致趋于完成 zeta 的 logarithmic-derivative germ。取 Cauchy 概率测度
\[
  \dd\mu(t)=\frac{\dd t}{\pi(1+t^2)},\qquad
  \cM=L^\infty(\R,\mu),\qquad \tau(f)=\int f\dd\mu.
\]
令 $P_{Y,\delta}(t)=\Re F_Y(\delta+it)$，$H_{Y,\delta}$ 为相应乘法算子。effects 正是 $0\leq\theta\leq1$ 的乘法算子，故
\[
 \inf_{0\leq\theta\leq1}\int P_{Y,\delta}\theta\dd\mu
 =-\int[P_{Y,\delta}]_-\dd\mu=-\frac{J_{Y,\delta}}\pi.
\]
此外，prime、continuum 与 Gamma 项可写成 signed unitary orbit Laplacians
\[
  H_{Y,\delta}=aI+L_- -L_+,
\]
其中每个 $L_\pm$ 都由长度半群的 edge squares 生成。因此“有限迹代数 + effect cone + signed Hodge squares”对 zeta 已无条件、显式存在。

\begin{theorem}[\statusT{} zeta 的精确存在性边界；预算蕴含 RH]
沿任一共尾日程 $Y\to\infty$、$\delta\downarrow0$，若
\[
  \sup\int_\R[P_{Y,\delta}(t)]_-\frac{\dd t}{1+t^2}<\infty,
\]
则 RH 成立。反之若 RH 假，则这一负指标沿 directed/cofinal 意义必发散。
\end{theorem}

第一部分是定理~\ref{thm:main} 的交换特例。反向若不成立，就能抽取一条有界负指标的共尾子列，再次推出 RH，矛盾。注意这不是 RH 的改写游戏：左侧的每个有限对象只用有限素数、Gamma 数据和显式正 kernel 构造；开放问题是证明其统一界。

\subsection{exact passive realization}

对完成函数 $\Phi$，乘法 current
\[
  H_x=M_{\Re(\Phi'/\Phi)(x+it)},\qquad x>0,
\]
对每个 $x$ 都非负，当且仅当 $\Phi'/\Phi$ 是右半平面的 positive-real 函数；这又等价于全部零点在中心线上及 Pick resolvent 的正实现。故\emph{精确正}结构的存在性与 RH/GRH 等价，不能作为无条件输入。本文的 bounded-index 定理有意把目标降为近似 current 的有限负深度。

\section{无条件有限化与等价判据族}

笔记 004--143 从不同坐标系反复检验同一 Hodge 内容。它们的共同成果不是许多个彼此独立的“RH 证明”，而是以下可交换的有限化字典。

\subsection{prolate 与有限自伴截断}

有限 Weil 矩阵、prolate near-radical、residual 证书、谱隙和 Feshbach 消元给出一条严格桥：若有限自伴行列式在临界带内局部一致收敛到 $\Xi(z)=\xi(1/2+iz)$，则 Hurwitz/Rouch\'e 保持零点实性，因而推出 RH。相关 prolate 与 zeta spectral-triple 构造见\cite{CC2021,CCMProlate,CCM2025}。仅有有限低零点的高精度逼近、$L^2$ 向量接近或极小 residual 都不够；需要 simple-even 最低态、统一谱隙以及足以控制 Mellin 变换的指数加权 $L^1$ 收敛。

\subsection{prime graph、Green kernel 与 wavelet}

prime translations 可写成正 graph squares；Sturm--Liouville、Bessel、Riesz、biharmonic、Sobolev、wavelet 与 annular frames 给出大量正 Gram。Tate/continuum 主模态可由 primitive dilation polynomial 消去。经过 Poisson sampling 和 Taylor 压缩，每个尺度的有效 rank 可降到
\[
  O\!\left(\frac{\log X}{\log\log X}\right),
\]
甚至对 off-center detection 可用 rank-one annular coordinate。所有这些有限 forms 对任意 coefficient vector 都正；仍未知的是实际 $\Lambda-1$ 向量的 Hodge 范数是否为 $X^{o(1)}$。

\subsection{Nyman--Beurling 与容量分支}

divisibility incidence、M\"obius Green form、Nyman--Beurling finite distance、Farey 边界、cotangent/Kloosterman 接口和 response capacity 形成另一条有限 Hodge 链；其解析基准见\cite{NymanBeurling,BaezDuarte,Burnol}。它精确恢复若干 RH 等价距离，并产生可计算的 Schur-short、cell-wedge 与 Vandermonde coercivity。但 positivity-only、sparse local frames、acyclic completion 和 response-zero 压缩均出现严格 no-go：若缺少 collective arithmetic cancellation，就不能从局部强制性推出中心线。

\subsection{passivity 与 capped cone}

positive-real 阻抗、Cayley inner semigroup、Bochner cone、capped Loewner SDP 和 stationary quadrature 最终汇入第 5 节的 Cauchy finite-trace current。这个分支的重要修正是：所需结论不是逐 character 的 Loewner 下界，而是负井的“深度乘容量”总和有界。

\section{平方根核心与调制 Type I/II Hodge 结构}

\subsection{layer cake 与无条件 localization}

对可测自伴 symbol $H$，
\[
  \int H_-\dd\mu=\int_0^\infty\mu\{H<-u\}\dd u.
\]
若在频率块 $I$ 上
$H=A+R+\Re Q$，$A\geq b>0$、$|R|\leq r<b$，令 $\beta=b-r$，则对 $p>1$，
\[
  \int_IH_-\dd t
  \leq \frac{(p-1)^{p-1}}{p^p}\frac{\int_I|Q|^p\dd t}{\beta^{p-1}}.
\]
\begin{theorem}[\statusT{} square-root Cauchy-core localization]\label{thm:sqrtcore}
固定 $0<\alpha<1/3$，令
\[
 \delta_Y=(\log Y)^{-\alpha},\qquad
 T_0(Y)=Y^{1/2}\delta_Y^{-3/2}\log Y.
\]
则上述 zeta Abel candidates 满足
\[
 \int_{|t|\geq T_0(Y)}[-\Re F_Y(\delta_Y+it)]_+
 \frac{\dd t}{1+t^2}=o(1).
\]
\end{theorem}

\begin{proof}
取固定 $A>2$，令 $N=\lceil AY\log Y\rceil$、$T_1=Y(\log Y)^6$，把 $S_Y$ 截成长度 $N$ 的 Dirichlet polynomial。指数 Abel tail 在 $T_0\leq|t|\leq2T_1$ 上为 $o(\log|t|)$；分部积分给 $I_Y(1/2+\delta_Y+it)=o(1)$，而 vertical Stirling 给
$\Re A_\infty(1/2+\delta_Y+it)\geq c\log|t|$。故每个 dyadic shell $T\leq|t|<2T$ 有 barrier $\beta_T\asymp\log T$。

写 $a_n=\Lambda(n)n^{-1/2-\delta_Y}e^{-n/Y}$。初等估计给
$\sum_{n\leq N}|a_n|^2\ll\delta_Y^{-3}$；Montgomery--Vaughan weighted Hilbert inequality\cite{MontgomeryVaughan} 给
\[
 \int_T^{2T}\left|\sum_{n\leq N}a_nn^{-it}\right|^2\dd t
 \ll\delta_Y^{-3}(T+Y\log Y).
\]
在 barrier inequality 中取 $p=2$，再乘 Cauchy 权 $\ll T^{-2}$，该 shell 的负指标至多
\[
 \ll\delta_Y^{-3}\left(\frac1{T\log T}
 +\frac{Y\log Y}{T^2\log T}\right).
\]
从 $T_0$ 到 $T_1$ 作 dyadic 求和，得到 $o(1)+O((\log Y)^{-2})$。在 $|t|\geq T_1$，相同 mean-square/Stirling 估计的绝对收敛高尾为
$O(\delta_Y^{-3}/(T_1\log T_1))=o(1)$。合并即得。
\end{proof}

因此 RH 的 bounded-defect 路线只剩 $|t|<T_0(Y)$ 内负指标的一致有界性。相比先前的 $Y\operatorname{polylog}Y$ 核心，这无条件缩短了近一个平方根。

\subsection{调制 Gallagher--Fej\'er 恒等式}

对 lag measure $\nu$，记
\[
  D(t)=\int e^{-it\lambda}\dd\nu(\lambda),\qquad
  \dd\nu_\tau=e^{-i\tau\lambda}\dd\nu(\lambda),\qquad
  g_{\tau,h}(x)=\nu_\tau([x,x+h]).
\]
则有精确 Plancherel 恒等式
\[
 \int_\R|g_{\tau,h}(x)|^2\dd x
 =\frac1{2\pi}\int_\R|D(\tau+s)|^2
 \frac{4\sin^2(sh/2)}{s^2}\dd s.
\]
原点附近的频率平移必须同时在 lag measure 上产生相位 $e^{-i\tau\lambda}$；这是 Gallagher 大筛引理\cite{Gallagher}在移动频率块中的正确形式。删除该 modulation 的普通 Selberg energy 无法统一控制所有高度块。

取 $\cH_h=L^2(\R)$、$e_\lambda=\1_{[\lambda-h,\lambda]}$，则
\[
 \langle e_\lambda,e_\mu\rangle=(h-|\lambda-\mu|)_+,
 \qquad e_{\log(dm)}=U_{\log d}e_{\log m}.
\]
于是 Dirichlet convolution 精确变成 logarithmic translation tensor synthesis；Vaughan/Heath--Brown Type I/II 分解因此成为 Hodge 结构的一部分，而非外加估计技巧。

\section{Vaughan 分解、canonical quotient 与矩阵通道}

\subsection{精确 centered Vaughan 恒等式}

令 $\mu_1=\mu\1_{n\leq U}$、$\mu_2=\mu-\mu_1$，$\Lambda_1=\Lambda\1_{n\leq V}$、$\Lambda_2=\Lambda-\Lambda_1$。在 Vaughan 恒等式的原始 Type I/II 语境\cite{Vaughan}下，本文使用的 centered 版本由 $\log=\Lambda*1$ 与 $\mu*1=\varepsilon$ 逐行推出；它逐整数精确满足
\[
\boxed{
\Lambda-1=(\mu_1*\log-1)-(\mu_1*\Lambda_1*1)
 +(\mu_2*\Lambda_2*1)+\Lambda_1.}
\]
连续背景可在四个 component 间作 affine gauge 分配；有限组 cutoffs 还可在 simplex 上联合优化。完整 physical vector 与 gauge 无关，component-diagonal majorant 随 gauge 改变。有限实验显示优化收益只有几个百分点，说明仅调 cutoff 不会产生数量级突破。

\subsection{Canonical 二通道 quotient 与算术长度局部化}

令
\[
 a_U=\mu_{\leq U}*1,\qquad X=a_U*\Lambda_{>V}.
\]
四个未中心化分量在 physical synthesis 下有一个不依赖数据拟合的精确商：
\[
 I_{U,V}=P_1+P_2+P_4=X+\Lambda_{\leq V},\qquad
 II_{U,V}=P_3=\mu_{>U}*\Lambda_{>V}*1.
\]

\begin{theorem}[\statusT{} canonical two-channel Vaughan quotient]\label{thm:two-vaughan}
逐整数有
\[
 \Lambda=I_{U,V}+II_{U,V}.
\]
而且
\[
 II_{U,V}(n)=0,\qquad
 I_{U,V}(n)=\Lambda(n)
 \quad\bigl(n<(U+1)(V+1)\bigr).
\]
若四分量 Hilbert Gram 为 $G_4$，置
\[
 S=\begin{pmatrix}1&1&0&1\\0&0&1&0\end{pmatrix},
 \qquad G_2=SG_4S^*,
\]
则 $G_2\geq0$ 且
\[
 (1,1)G_2(1,1)^*=\1_4^*G_4\1_4.
\]
因此 physical energy 可无误差地写成 canonical $2\times2$ Type I/II Gram；这里没有 principal-angle、tail 或 Feshbach approximation。
\end{theorem}

\begin{proof}
由 $\log=\Lambda*1$，$P_1+P_2=\mu_{\leq U}*\Lambda_{>V}*1=X$。再由 $\mu*1=\varepsilon$，
\[
 P_3=(\mu-\mu_{\leq U})*\Lambda_{>V}*1
    =\Lambda-\Lambda_{\leq V}-X.
\]
这给分解。$P_3$ 的每个非零三因子项满足 $n=dmr$、$d>U$、$m>V$、$r\geq1$，故有 support gap；Type I agreement 随即来自总和恒等式。Gram 结论由 $S^*(1,1)^*=\1_4$。
\end{proof}

support gap 同时是一条 cutoff 守恒律：把 Type II 推到更大的 $UV$，会让 Type I 在完全相同的初始区间逐项复制 $\Lambda$，所以 ``small Type II'' 不能自动控制目标。不过这一区间的低算术长度部分可以无条件删除。令
\[
 L_U=\sum_{n\leq U}\Lambda(n)n^{-\sigma}e^{-n/Y}n^{-i\tau}e_{\log n},
 \qquad \tfrac12\leq\sigma\leq\tfrac34.
\]
在所用 cofinal schedule 中 $\sigma=1/2+\delta_Y\to1/2$，故该范围最终自动满足；有限个初始尺度吸收到常数。
由 $\psi(x)\ll x$、partial summation 与 $\|e_{\log n}\|=\sqrt h$，
\[
 \|L_U\|^2\ll hU^{2(1-\sigma)}\leq hU.           \tag{*}
\]

\begin{corollary}[\statusT{} arithmetic-length localization]\label{cor:length-local}
在 dyadic height block $T_k=2^k$ 上有 $h_k\asymp T_k^{-1}$、barrier $\beta_k\gg\log(e+T_k)$。对任意固定 $A>0$ 取
\[
 U_k=\max\!\left(1,\left\lfloor
 \frac{T_k}{\log^A(e+T_k)}\right\rfloor\right).
\]
则一致于 Abel scale、水平偏移和 modulation center，
\[
 \sum_k\frac{\|L_{U_k}\|^2}{\beta_k}<\infty.
\]
写 $p_{I,k}=L_{U_k}+r_{I,k}$，则 $r_{I,k}$ 的系数支撑于 $n>U_k$，$p_{II,k}$ 支撑于 $n\geq(U_k+1)(V_k+1)$。因此若
\[
 \sup_Y\sum_{k\in\mathrm{core}}
 \frac{\|r_{I,k}+p_{II,k}-c_k\|^2}{\beta_k}<\infty,              \tag{**}
\]
其中 $c_k$ 是同一 block 的 continuum vector，则 RH 成立。
\end{corollary}

\begin{proof}
由 (*)，第一个级数至多常数乘
$\sum_k\log^{-(A+1)}(e+2^k)<\infty$。定理~\ref{thm:two-vaughan} 给两个 support assertions。完整 block vector 为 $L_{U_k}+r_{I,k}+p_{II,k}-c_k$；二向量 Cauchy--Schwarz、(**) 与外部平方根核心的无条件估计给调制 Gallagher budget，故 bounded finite-trace Hodge--Weil 定理推出 RH。
\end{proof}

该推论把尚未证明的量从整个 prime current 缩到 $n>T/\log^A T$ 的 signed truncated-M\"obius Type I residual、support 更晚的 Type II residual及其精确 cross Gram。它仍是 RH 强度输入，但不再含每块的低算术长度。
\subsection{矩阵 Bessel 与 Feshbach--Loewner}

令四个 component vectors 的 Gram 为 $G\geq0$，physical direction 为 $e=(1,1,1,1)^T$，则真实能量是 $e^*Ge$。若 $G\leq M$ 且
\[
  \sup_Y\sum_k\frac{e^*M_{Y,k}e}{\beta_{Y,k}}<\infty,
\]
则由调制能量判据和平方根核心 localization 推出 RH。

对一个固定 rank-$q$ 投影 $P$，写
\[
  G=\begin{pmatrix}A&B\\B^*&C\end{pmatrix}.
\]
任取 $\epsilon>\|C\|$，有严格 Loewner majorant
\[
 G\leq
 \begin{pmatrix}
 A+B(\epsilon I-C)^{-1}B^*&0\\0&\epsilon I
 \end{pmatrix}.
\]
这是 upper Feshbach completion；它允许公共通道不逐块对角化，并把 core--tail coupling 的代价保留为正 resolvent correction。

\subsection{Laurent 审计与 incidence 基}

四个未中心化 Vaughan Dirichlet series在 $s=1$ 的双极点与单极点向量为
\[
 d=m_0(1,0,-1,0),
\quad
 r=(m_1,-m_0\ell_0,1+m_0\ell_0-m_1,0).
\]
平衡 continuum residue 后，两条 singular directions 都落在 $e^\perp$。因此完整 Laurent singular plane 不可能同时是包含 physical direction 的二维 core；这是严格的维数障碍，而非数值不稳定。

使用正交 incidence 基
\[
 e_0=\frac{(1,1,1,1)}2,\quad
 u=\frac{(1,-1,0,0)}{\sqrt2},\quad
 v=\frac{(0,0,-1,1)}{\sqrt2},\quad
 c_0=\frac{(1,1,-1,-1)}2,
\]
可选冻结 tilted core
\[
 \cC=\operatorname{span}\{e_0+\tfrac1{20}c_0,\ 3u+v\}.
\]
其整数 core seeds 和 tail seeds 可取
\[
 (21,21,19,19),\ (3,-3,-1,1);\qquad
 (19,19,-21,-21),\ (-1,1,-3,3).
\]
physical vector 在 core/tail 的质量比例精确为 $400/401$ 与 $1/401$。这一基由 component labels 固定，不依赖零点和高度。

\subsection{冻结通道的精确 Type I/II 公式}

记
\[
 a_U=\mu_{\leq U}*1,\quad X=a_U*\Lambda_{>V},\quad
 Y=a_U*\Lambda_{\leq V},\quad V_0=\Lambda_{\leq V},\quad L=\Lambda.
\]
冻结正交变换后的四个通道满足
\begin{align*}
 Z_0&=\frac{2X+19L}{2\sqrt{401}},&
 Z_1&=\frac{4X+6Y-L+2V_0}{2\sqrt5},\\
 Z_2&=\frac{40X-21L}{2\sqrt{401}},&
 Z_3&=\frac{2X-2Y-3L+6V_0}{2\sqrt5},
\end{align*}
并且
\[
  L=\frac{40}{\sqrt{401}}Z_0-\frac2{\sqrt{401}}Z_2.
\]
这把 tail 明确化为两个可供 Type I/II large sieve 处理的 residual $Z_2,Z_3$。但 $Z_0,Z_1$ 仍显式含 $L$，故“证明 tail 很小”不能自动控制目标。

\begin{theorem}[\statusT{} component compression no-free-lunch]\label{thm:nofree}
令 $P$ 为 component label space 上任意投影，$Q=I-P$，$e$ 为 physical vector。若 $Pe\neq0$，则对每个 $M>0$ 存在 rank-one $G_M\geq0$，使
\[
  PG_MQ=0,\qquad QG_MQ=0,\qquad
  \langle G_Me,e\rangle=M\|Pe\|^2.
\]
若 $Pe=0$，则可把任意大物理能量完全放入 tail，同时令 core 与 coupling 为零。
\end{theorem}

\begin{proof}
当 $Pe\neq0$ 时取 $u=Pe/\|Pe\|$、$G_M=M|u\rangle\langle u|$。第二种情形取 $u=e/\|e\|$。
\end{proof}

定理~\ref{thm:nofree} 是当前最重要的反面结论。低秩、稳定 principal angles、小 Feshbach correction 和小 tail edge 都不能替代 core 估计。具体 zeta Grams 当然未必允许任意 rank-one amplification；真正的证明必须利用它们来自 primes、Gamma、continuum 和 convolution identities 的事实，构造一个禁止这种放大的数域 Hodge--Riemann 关系。

\input{paper-sections/mobius-threshold-hodge}
\input{paper-sections/profinite-diagonal-evacuation}
\input{paper-sections/selberg-volterra-cross-fiber}
\input{paper-sections/selberg-profile-equivalence}
\input{paper-sections/capped-correspondence-kernel}
\input{paper-sections/reciprocal-barrier-schur}
\input{paper-sections/square-completed-schur-l4}
\input{paper-sections/tracial-lattice-clipping}

\section{存在性审计：已经构造了什么，尚缺什么}

\begingroup\small
\begin{longtable}{@{}>{\raggedright\arraybackslash}p{0.16\textwidth}p{0.16\textwidth}p{0.14\textwidth}p{0.42\textwidth}@{}}
\toprule
结构环节 & 有限域 & 经典 zeta & 当前状态与缺口\\
\midrule
\endfirsthead
\toprule
结构环节 & 有限域 & 经典 zeta & 当前状态与缺口\\
\midrule
\endhead
全局谱/上同调载体 & \(\ell\)-adic/晶体上同调 & 多种 finite Gram、GNS、Abel 与 tracial carriers & finite carriers 无条件；canonical global cohomology 未构造\\
迹/行列式公式 & Lefschetz 迹公式 & Weil 显式公式、Hadamard/正规化 determinant & 解析 divisor 接口存在；需与最终极限结构严格相容\\
中心对称 & Poincar\'e 对偶 & 完成函数的函数方程 & 无条件\\
primitive/Tate 分解 & Hard Lefschetz primitive decomposition & continuum subtraction 与 dilation polynomial & finite/filtered 形式无条件\\
正极化 & Hodge--Riemann/Rosati & positive Green、wavelet、triangular、moment Grams & 载体正性无条件，但不等于实际算术 tightness\\
精确纯性 & Deligne/Weil & exact passive/adjoint realization & 对 zeta 与 RH 等价，未证\\
filtered carrier & 自动由精确极化 & Euler、Gram、dilation、定性 visibility 与 tails & adaptive Sobolev 对偶分离无条件；任意 realization 仍需另证\\
算术 tightness & Hodge--Riemann 已给 & 实际向量 subpower norm & 与 RH 等强，未证\\
有限迹 Hodge 指数 & 负指标为零 & Cauchy effect cone 与 signed orbit current & 代数和 current 无条件；统一负迹界未证\\
有上界 correspondence kernel & 有效代数对应/极化子空间 & $r,\kappa-r$ 双正定与 joint orbit functional & ambient kernel、精确对偶与 block gluing无条件；统一 index界未证\\
reciprocal-barrier shorting & Hodge 正背景与 primitive projection & $d\mu/B$ kernel、harmonic predictor与 Schur residual & 变分式和有限 Gram无条件；amplitude路线由下一行的 quadratic lift替代\\
square-completed background & quadratic Hodge lift & Schur residual与 fourth correspondence Gram & lift及sharpness无条件；square-root fourth-moment budget未证\\
tracial lattice clipping & order functional calculus & operator barrier与 clipped residual & finite-trace theorem无条件；zeta单侧 sublevel budget未证\\
平方根核心 & 不需要 & exterior defect 为 \(o(1)\) & 无条件；core bound 未证\\
Type I/II tensor & correspondence calculus & centered Vaughan + Fej\'er incidence & 精确、无条件\\
canonical 二通道 & physical quotient & exact $G_2=SG_4S^*$ 与 support-product gap & quotient及 $n\leq T/\log^A T$ 消去无条件；$V$ 仅为 hard-channel gauge\\
M\"obius threshold Hodge & divisor threshold complex & harmonic/heat supertrace、boundary shell 与 gcd Gram & 有限结构无条件；与 interval incidence 的全局 finite-index compatibility 未证\\
profinite gcd polarization & divisibility cylinders/ideal lattice & twists与 ideals上的 positive Gram；单调加权 second moment & hard diagonals在 $T\geq\log^K Y$ 无条件消去；off-diagonal 未证\\
Selberg--Volterra/profile wall & 前缀场 & exact Abel transport、microscopic incidence与 fixed-dilation Mellin transform & transport及 $N\leq T^{2-\eta}$ 消去无条件；polylog full profile与 RH 等价\\
经验低秩通道 & 权重分解 & frozen tilted rank two & 仅作 finite diagnostic；不再是必要 proof input\\
core Hodge relation & Hodge 指数定理 & 尚无对应物 & 当前决定性缺口\\
\bottomrule
\end{longtable}
\endgroup

对 primitive Dirichlet $L$ 函数，可将非自对偶对象与其对偶成对，局部相位进入 twisted orbit edges；Gamma 因子和 ramified 项成为有限或 regularized corrections。对固定次数 automorphic 数据，还需逐例验证 Rankin--Selberg diagonal、local coefficient growth 和 conductor-uniform 误差。若 local temperedness 本身是 Ramanujan 型猜想，就不能把它伪装成无条件输入。

Deninger 的叶状动力系统纲领与 Connes 的 adele 类空间迹公式\cite{Deninger1998,ConnesTrace}已经提供了严肃的全局几何候选，但现有结果尚未同时给出本文所需的 determinant、全局 trace 和非循环极化。它们应被视为可能产生 core Hodge relation 的来源，而不是已完成的 G1--G4 包。

\section{主要 no-go 与逻辑防线}

研究过程中得到的反面结论与正面定理同样重要：
\begin{enumerate}
  \item \textbf{函数方程不够。}它只给四元组对称，不能给中心线。
  \item \textbf{逐素数纯性不够。}全局零点是跨所有素数解析耦合后的对象。
  \item \textbf{抽象 Hilbert--P\'olya 存在性会循环。}若谱被预先指定为零点，则 adjoint 关系本身已经含 RH。
  \item \textbf{有限数值零点不够。}需要临界带内局部一致的整函数极限。
  \item \textbf{有限正 Gram 不够。}正性对所有 coefficient vectors 成立，但实际算术向量的临界大小仍未知。
  \item \textbf{ordinary short-interval energy 不够。}高度平移后必须保留 $e^{-i\tau\lambda}$ modulation。
  \item \textbf{pointwise Loewner 下界过强。}bounded negative trace 允许深而窄的负井。
  \item \textbf{稀疏/local coercivity 不够。}Nyman/Farey 分支需要 collective frame 或外部周期信息。
  \item \textbf{Laurent singular plane 不等于 physical core。}前者位于 $e^\perp$，两维不可能同时容纳全部 singular directions 和 $e$。
  \item \textbf{tail/coupling-only 压缩不够。}定理~\ref{thm:nofree} 要求独立 core Hodge 输入。
  \item \textbf{完整 hard cross Gram 不是较弱目标。}Hadamard 正规形把所需 cross cancellation 加 primitive energy 精确还原为 physical energy。
  \item \textbf{raw factorization Bessel 不一致。}exact-product fibers 的 multiplicity 无界；必须先沿 fibers 保留 M\"obius signs 作商化。
  \item \textbf{local Hodge direct sum 不产生 cross maps。}fiberwise McKean--Singer只给 coefficient lift；不同 divisor complexes间的 physical Gram entries仍需 arithmetic correspondences。
  \item \textbf{scalar second moment 不控制 prefixes。}长度 $H$ 的 constant block有 coefficient mass $H$ 而 terminal prefix square为 $H^2$；profinite diagonal不能自动推出 Selberg profile。
\end{enumerate}

这些 no-go 共同保证：本文的“广义结构定理”没有通过重新命名 RH 来制造证明。每条充分条件都清楚标出仍需证明的算术量。

\section{有限实验：能说明什么，不能说明什么}

实验在方法选择上提供了三项稳定信息：
\begin{enumerate}
  \item Cauchy 最深负井可随参数加深，而 sampled negative capacity 反而下降，支持 finite-trace 而非 operator-norm 目标；
  \item 四 component Vaughan Grams 的 physical energy 在小样本中几乎由两个通道承载，公共二维 Feshbach upper excess 曾低于 $0.74\%$，冻结整数基的若干 held-out blocks excess 低于 $0.03\%$；
  \item exact incidence/tilted 基把物理质量按 $400/401$ 与 $1/401$ 分到 core/tail，并导出显式 Type I/II residual。
\end{enumerate}

定理~\ref{thm:two-vaughan} 的二通道 quotient 是逐系数代数恒等式，不依赖这些实验；它使上述经验 bases 降为诊断坐标，而非下一证明的必要结构。

这些数字全部是 finite、floating-point、midpoint/truncated diagnostics；尚未同时包含 interval arithmetic、无限 continuum tail、prime tail、全部 dyadic heights 和共尾极限。因此它们不能证明 bounded $J_Y$，更不能作为 RH 的统计证据。它们的合法用途是发现 invariants、选择 proof norm、构造反例和优先排序下一步估计。

\section{下一步研究议程}

文档 170--173 已依次把 effect cone压缩为长度侧结构、把 joint index转成 reciprocal-barrier Schur residual、以 square completion删除 pointwise amplitude，并由 quartic-capacity审计转向 finite-trace lattice clipping；当前不应再平行堆叠等价坐标或只优化有限 cutoffs。决定性任务是估计 explicit predictor越过负 barrier的 one-sided clipped current。建议按以下顺序推进。

\subsection{第一阶段：threshold-Hodge superconnection current}

固定推论~\ref{cor:length-local} 的 $U_T=T/\log^A T$。profinite polarization已经无条件控制 $T\geq\log^K Y$ 的 arithmetic coefficient diagonals；Selberg--Volterra transport又给出 ordinary prefix field 到 modulated current 的显式 cross-fiber map。现在应将它与定理~\ref{thm:mobius-hodge} 的 reduced complexes、Hodge Laplacians及 profinite cylinders组合成自伴矩阵值 superconnection current $\mathbb H_{Y,\delta}(t)$，要求：
\begin{enumerate}
  \item scalar Abel zero obstruction 由其负谱迹控制；
  \item effect cone 以定理~\ref{thm:capped-correspondence} 的双正定 kernels实现，并以定理~\ref{thm:lattice-clip} 的 finite-trace clipped background控制其负指标；
  \item non-harmonic divisor faces由 $\partial+\partial^*$ supersymmetry配对消去，负方向只读取 harmonic boundary classes；
  \item 不把与 RH 等价的 full $\mathcal J_N(s)\ll Ns\,\operatorname{polylog}N$ 当作较弱输入；改为在 square-root resonance wedge构造只读取 barrier 负方向的 signed prefix/layer-cake版本；
  \item 结构由 prime/Gamma orbit correspondences 与 divisor incidence构造，不引用零点位置。
\end{enumerate}

\subsection{第二阶段：非循环 core 不等式}

需要证明一种排除 rank-one core amplification 的算术关系。可接受的候选形式包括：
\begin{enumerate}
  \item indefinite intersection form 加有限 Hodge index；
  \item prime/Gamma-built positive current 对负谱部的直接 majorization；
  \item 利用 M\"obius 与 Type I/II 结构的 signed orbit comparison；
  \item 一个由 arithmetic dynamics 推出的双向 tempered 或 adjoint/similitude 关系。
\end{enumerate}
任何只估计完整 $\|I_{U,V}\|$ 而不利用 continuum cancellation或 signed M\"obius structure 的方案，都必须面对定理~\ref{thm:two-vaughan} 的 target-agreement range，可能把 $\Lambda$ 原样放回右端。

\subsection{第三阶段：tail 与 coupling 的解析预算}

以 boundary-shell 公式 (BS) 压缩 $b_U$ 后，只需处理定理~\ref{thm:wedge-evacuation} 留下的 $N>T^{2-\eta}$ square-root resonance wedge。diagonal 已由 profinite gcd Gram处理，cross entries已由式 (SH) 归约；但定理~\ref{thm:profile-rh} 表明不能再以证明 full (SP) 作为“较弱中间目标”，而式 (QC) 又表明 unscaled fourth budget可能过强。新的解析目标是按定理~\ref{thm:lattice-clip} 构造 threshold-complex/Type I--II joint predictor $C_T$，并证明
\[
 \sum_TK_T(\varepsilon)<\infty,\qquad
 K_T(\varepsilon)=\int_{I_T}
 \frac{(R_T-C_T^{[\varepsilon]})^2}{B_{0,T}}\dd\mu .
\]
低 polylog-height core再单独加入。应优先估计 bad set
$\{C_T<-(1-\varepsilon)B_{0,T}\}$ 的单侧 capacity；Schur--$L^4$ 仍可逐块作为比较证书，而非默认主目标。

Guth--Maynard 的 Dirichlet polynomial large-value bounds\cite{GuthMaynard}是 layer-cake路线的潜在输入之一，但其 theorem控制 bounded coefficients在 $1$-separated points上的 superlevel cardinality，并不自动给出 continuous balanced $L^2$ profile (SP)；coefficient normalization、maximal transfer、prime--continuum balancing与 threshold integration必须单独证明。曾声称 log-power intervals almost-all estimate的 \texttt{arXiv:1009.6121} 已由作者因关键引理错误撤回，本文不把该 claim作为输入。

\subsection{第四阶段：严格有限证书与广义化}

用 interval arithmetic 和 exact Cauchy cell masses替代浮点 midpoint 审计，给出共尾 schedule 上的可验证上下界。若 zeta 路线闭合，再把同一结构迁移到 paired Dirichlet 和固定次数 Gamma--Euler 数据，明确记录 degree、conductor、ramification 与 Rankin--Selberg 输入的统一性。

\section{结论}

从 Weil 猜想中真正可抽离的，不是“存在一个神秘上同调”这句话，而是一组可审计的机制：全局 trace compatibility、中心对偶、正极化或有限负指标，以及实际算术 cycle 的统一 Hodge 控制。本文在有限维、tempered 和有限迹层次给出严格结构蕴含，并把 filtered 层次修订为含定量 dual/Riesz separation 的条件框架。对 zeta，Mellin-coherent adaptive Sobolev carrier 的 separation 已由定理~\ref{thm:mellin-dual} 无条件建立；任意 Gram 的全空间 separation 仍非自动。真正保持 RH 强度的开放输入是实际算术 cycle 的统一 subpower Hodge 估计，平方根核心 localization 只缩小了该估计的范围，并未证明它。

最新的 canonical quotient 又把经验通道降为一个精确 Type I/II 对象，并无条件删除了 $n\leq T/\log^A T$ 的低算术长度；冻结 basis 现在只保留为诊断工具。hard-channel gauge 与 Hadamard no-go 进一步表明：完整 cross Gram 并不是较弱的新目标。截断 M\"obius defect 已被实现为 zeros-independent 的 threshold-complex Hodge heat supertrace；profinite divisibility polarization又把其 second moment转移到全部单调权，并无条件消去 $T\geq\log^K Y$ 的 hard coefficient diagonals。Selberg--Volterra identity提供真正的 cross-fiber transport；microscopic incidence又无条件删除 $N\leq T^{2-\eta}$ 的 wedge。与此同时，telescoping--Mellin theorem证明 polylogarithmic full profile本身与 RH 等价。capped-correspondence theorem把频谱 effects翻译成 $r$ 与 $\kappa-r$ 的双正定性；reciprocal-barrier theorem再把 joint index化为 harmonic predictor后的 Schur residual。square completion删除了 pointwise amplitude，但 quartic-capacity下界又揭示 unscaled fourth budget可能过强。finite-trace lattice clipping最终把这两个代价统一替换为 one-sided quadratic residual，而且允许 background与 current不对易。下一次真正的突破应在剩余 square-root resonance wedge内证明该 clipped residual可和。若由此得到 Cauchy 负指标一致有界，定理~\ref{thm:main} 将直接给出 RH；对一般 Gamma--Euler 数据，还必须验证相应 ideal/local-system counting inputs。

非构造路线现在也有了严格边界：定理~\ref{thm:four-moment-nogo} 排除任何只用
四个标量矩的“自动正性”，而定理~\ref{thm:localizer-completion} 证明增长的 mixed
arithmetic localizers 足以非构造地产生全局 tracial positive representation。
因此下一次结构性突破可以采取两种等价可审计形式：直接证明 clipped residual
可和，或证明一个 divisor-visible operator system 的有限 localizers 一致近正；
二者都不能由低阶标量矩或裸紧致性免费得到。

\appendix
\section{202 篇笔记的主题映射}

下表给出原始笔记到本文论证链的完整范围映射。编号 003 为文献与证据审计；
其余 201 篇为研究笔记。

\begingroup\small
\begin{longtable}{@{}p{0.09\textwidth}p{0.28\textwidth}p{0.51\textwidth}@{}}
\toprule
编号 & 主题 & 在本文中的角色\\
\midrule
\endfirsthead
\toprule
编号 & 主题 & 在本文中的角色\\
\midrule
\endhead
001--003 & PLF 主定理、数域存在性审计、文献边界 & 第 2--3 节与参考文献的逻辑基线\\
004--014 & prolate 最低态、有限 Weil 矩阵、residual、谱簇与 Ritz/Feshbach & 第 7 节的有限自伴截断桥及其收敛缺口\\
015--030 & filtered near-radical、prime graph、resonance、紧 Hodge 算子与 trace-class defect & FPW 的早期有限化和负谱/共振结构\\
031--043 & Mellin 边界、Besicovitch/Abel/Chebyshev、Green--Bessel、dyadic/windowed cores & 正 kernel、边界能量与低频核心\\
044--058 & Riesz、dilation/de Rham、wavelet、primitive、moment/Sobolev/sampling 与 FPW 条件框架 & 第 4 节的 filtered primitive Weil package\\
059--077 & annular detector/frame、Abel GNS、RKHS、prefix/dual Laplacian、tempered category、multiplicity 与 trace determinant & strong GNS、tempered 酉化和 tracial multiplicity\\
078--087 & Ihara、Schatten no-go、divisibility、Nyman--Beurling、M\"obius/Farey、zero jets 与 capacity & 第 7 节的图 zeta和 Nyman 有限 Hodge 路线\\
088--105 & parabolic Farey、cotangent/Kloosterman、response leverage、Schur/cell/Vandermonde 与 sparse recovery & collective-frame 必要性和局部 coercivity 边界\\
106--130 & susceptibility、unit-cell、rank update、Pad\'e、polyhedral certificates，以及 Mertens、cycles、purity、completion、anomaly & response-capacity 与 positive-only/no-free-purity 审计\\
131--143 & positive-real/passivity、Abel resonance、inner/Bochner cone、capped SDP、quadrature、cofinal discretization & 第 5--6 节的 Cauchy effect cone 和 fully finite 接口\\
144--150 & zero-wave obstruction、bounded-defect normality、quadratic no-go、signed orbit Laplacian、finite trace、layer cake & 定理~\ref{thm:main} 及平方根核心 localization\\
151--158 & modulated Gallagher、Fej\'er、Vaughan、background/cutoff gauge、matrix Bessel、Feshbach、frozen out-of-sample & 第 8--9 节的 Type I/II 矩阵通道\\
159--162 & Laurent obstruction、incidence cone、tilted frozen basis、精确通道与 compression no-free-lunch & 当前边界与第 11 节研究议程\\
163 & Mellin 横坐标、compact-frequency dual 与 Schur margin & 第 4 节定理~\ref{thm:mellin-dual} 的定量分离机制\\
164 & canonical 二通道 Vaughan quotient、support-product 守恒与 arithmetic-length localization & 定理~\ref{thm:two-vaughan} 与推论~\ref{cor:length-local}\\
165 & hard-channel cutoff gauge、Hadamard cross-Gram 循环性与 factorization-fiber 障碍 & 定理~\ref{thm:hard-gauge} 及 hard-channel no-go\\
166 & M\"obius threshold complex、Hodge heat supertrace、boundary shell 与 positive gcd Gram & 定理~\ref{thm:mobius-hodge} 及下一阶段 signed superconnection 接口\\
167 & length-scalar supertrace、profinite divisibility polarization 与 hard diagonal evacuation & twists/ideals的 global gcd Gram及 high-core diagonal localization\\
168 & Selberg--Volterra cross-fiber transport、multiscale prefix profile 与来源审计 & full near-product 的 Hardy--Selberg 条件接口\\
169 & microscopic prefix bound、平方根共振楔与 Selberg profile 的 RH 等价性 & 无条件 wedge localization及 full profile 的逻辑强度审计\\
170 & capped positive-definite correspondences、joint orbit index与 block gluing & 长度侧 generalized Weil structure及 zeta存在性边界\\
171 & reciprocal-barrier kernel、harmonic Schur shorting与 amplitude证书 & one-sided square-root target及 finite LMI接口\\
172 & square-completed positive background、Schur--$L^4$ shorting与 fourth correspondence Gram & amplitude-free bounded-index criterion及新 square-root moment target\\
173 & quartic-capacity审计、operator reciprocal barrier与 relative lattice clipping & noncommuting finite-trace structure及 one-sided zeta target\\
174--180 & 非构造存在性、finite cone separators、Hodge transgression、随机 Fej\'er、packet criterion与 operator-system repair & 增长 localizer 路线的有限凸几何、换位子与 capture 边界\\
181--196 & modulated complement、bicommutant、finite-word moments、soft Cauchy effects、Brownian primitive 与 response-specific Vaughan Gram & mixed arithmetic word systems、非构造 response compactness及其 no-go\\
197--202 & 部分 Weil 比例与深度、quadratic fourth channels、smooth-window cycle、Schatten boundary、product/ratio Gram 与 localizer completion & 独立比例论文及本节的 scalar-four-moment no-go、finite-satisfiability completion\\
\bottomrule
\end{longtable}
\endgroup

\section{核心定理的依赖关系}

\begin{center}
\begin{tabular}{c}
形式结构：Lefschetz + 相容 primitive 极化 $\Longrightarrow$ PLF similitude $\Longrightarrow$ pure weights\\[4pt]
$\Downarrow$ 抽象化\\[4pt]
双向 tempered dynamics $\Longrightarrow$ spectral radius purity\\[4pt]
$\Downarrow$ cyclic/filtered 化\\[4pt]
explicit formula + visible divisor modes + FPW tightness $\Longrightarrow$ center line\\[4pt]
$\Downarrow$ 非构造补全\\[4pt]
finite mixed localizers + Archimedean bounds $\Longrightarrow$ tracial positive representation\\[4pt]
$\Downarrow$ 数域弱化\\[4pt]
Cauchy effect realization + bounded finite-trace negative index\\
$\Longrightarrow$ normal family $\Longrightarrow$ logarithmic derivative 无右半平面极点\\
$\Longrightarrow$ self-duality $\Longrightarrow$ center line.
\end{tabular}
\end{center}

在 zeta 上，前述箭头左侧的解析载体均已构造；最后未闭合的输入是
\[
 \sup_{Y,\delta}\tau\bigl((H_{Y,\delta})_-\bigr)<\infty,
\]
或任何经本文定理严格推出它的非循环 core Hodge 估计。

\begin{thebibliography}{99}

\bibitem{Deligne1974}
P. Deligne, \emph{La conjecture de Weil I}, Publ. Math. IH\'ES \textbf{43} (1974), 273--307.
\url{https://publications.ias.edu/node/368}

\bibitem{GrothendieckStandard}
A. Grothendieck, \emph{Standard Conjectures on Algebraic Cycles}, Bombay Colloquium (1968/1969).
\url{https://webusers.imj-prg.fr/~leila.schneps/grothendieckcircle/StandardConjs.pdf}

\bibitem{KleimanWeil}
S. L. Kleiman, \emph{Algebraic Cycles and the Weil Conjectures}.
\url{https://agrothendieck.github.io/divers/kleimanweil.pdf}

\bibitem{WeilExplicit}
A. Weil, \emph{Sur les ``formules explicites'' de la th\'eorie des nombres premiers} (1952).
\url{https://cds.cern.ch/record/471308}

\bibitem{Deninger1998}
C. Deninger, \emph{Some analogies between number theory and dynamical systems on foliated spaces}, ICM 1998.
\url{https://ems.press/books/dms/246/4682}

\bibitem{ConnesTrace}
A. Connes, \emph{Trace formula in noncommutative geometry and the zeros of the Riemann zeta function}.
\url{https://arxiv.org/abs/math/9811068}

\bibitem{CCM2025}
A. Connes, C. Consani and H. Moscovici, \emph{Zeta Spectral Triples} (2025 preprint).
\url{https://arxiv.org/abs/2511.22755}

\bibitem{CC2021}
A. Connes and C. Consani, \emph{Spectral Triples and Zeta-Cycles}.
\url{https://arxiv.org/abs/2106.01715}

\bibitem{CCMProlate}
A. Connes, C. Consani and H. Moscovici, \emph{Zeta Zeros and Prolate Wave Operators}.
\url{https://arxiv.org/abs/2310.18423}

\bibitem{Titchmarsh}
E. C. Titchmarsh, revised by D. R. Heath-Brown, \emph{The Theory of the Riemann Zeta-Function}, 2nd ed., Oxford University Press, 1986.

\bibitem{NymanBeurling}
M. Balazard and E. Saias, \emph{The Nyman--Beurling Equivalent Form for the Riemann Hypothesis} (1998).
\url{https://www.esi.ac.at/preprints/esi623.pdf}

\bibitem{BaezDuarte}
L. B\'aez-Duarte, \emph{A Strengthening of the Nyman--Beurling Criterion for the Riemann Hypothesis}, Rend. Lincei \textbf{14} (2003), 5--11.

\bibitem{Burnol}
J.-F. Burnol, \emph{A Lower Bound in an Approximation Problem Involving the Zeros of the Riemann Zeta Function}, Adv. Math. \textbf{170} (2002), 56--70.

\bibitem{MontgomeryVaughan}
H. L. Montgomery and R. C. Vaughan, \emph{Hilbert's Inequality}, J. London Math. Soc. (2) \textbf{8} (1974), 73--82.

\bibitem{Gallagher}
P. X. Gallagher, \emph{A Large Sieve Density Estimate Near $\sigma=1$}, Invent. Math. \textbf{11} (1970), 329--339.

\bibitem{Vaughan}
R. C. Vaughan, \emph{Sommes trigonom\'etriques sur les nombres premiers}, C. R. Acad. Sci. Paris S\'er. A \textbf{285} (1977), 981--983.

\bibitem{SaffariVaughan}
B. Saffari and R. C. Vaughan, \emph{On the Fractional Parts of $x/n$ and Related Sequences. II}, Ann. Inst. Fourier \textbf{27} no. 2 (1977), 1--30.

\bibitem{GuthMaynard}
L. Guth and J. Maynard, \emph{New Large Value Estimates for Dirichlet Polynomials}, Ann. of Math. \textbf{203} (2026), 623--675.

\bibitem{Bass}
H. Bass, \emph{The Ihara--Selberg Zeta Function of a Tree Lattice}, Int. J. Math. \textbf{3} (1992), 717--797.

\bibitem{Simon}
B. Simon, \emph{Trace Ideals and Their Applications}, 2nd ed., AMS, 2005.

\bibitem{BurgdorfKlep}
S. Burgdorf and I. Klep,
\emph{The truncated tracial moment problem},
J. Operator Theory \textbf{68} (2012), 141--163.
\url{https://arxiv.org/abs/1001.3679}

\bibitem{MourrainSchmudgen}
B. Mourrain and K. Schm\"udgen,
\emph{Flat extensions in $*$-algebras},
Proc. Amer. Math. Soc. \textbf{144} (2016), 4873--4885.
\url{https://arxiv.org/abs/1406.4975}

\bibitem{DLMF}
NIST Digital Library of Mathematical Functions, §§5.9, 5.11.
\url{https://dlmf.nist.gov/}

\end{thebibliography}

\end{document}
